\documentclass[11pt]{article}
\usepackage[a4paper,margin=18mm]{geometry}
\usepackage{fontspec}
\setmainfont{Latin Modern Roman}
\usepackage{amsmath,amssymb,amsthm,amscd,graphicx,color}
\usepackage{xurl}
\usepackage[unicode,hidelinks]{hyperref}
\allowdisplaybreaks
\setlength\emergencystretch{3em}
\numberwithin{equation}{section}

\newtheorem{Theorem}{Theorem}[section]
\newtheorem{Corollary}[Theorem]{Corollary}
\newtheorem{Lemma}[Theorem]{Lemma}
 { \theoremstyle{definition}
\newtheorem{Definition}[Theorem]{Definition}
\newtheorem{Example}[Theorem]{Example} }


\begin{document}
\begin{center}\Large Integral-operator characterization of spectral functions of real Jacobi-type five-diagonal pencils\end{center}
\noindent\textbf{Stable ID:} 1887\par
\noindent\textbf{Authors:} Sergey M. Zagorodnyuk\par
\noindent\textbf{Original work:} The Inverse Spectral Problem for Jacobi-Type Pencils\par
\noindent\textbf{Version:} Official SIGMA 13 (2017) article 085, DOI 10.3842/SIGMA.2017.085; journal PDF and native TeX acquired 2026-10-01, exact bytes pinned by SHA256\par
\noindent\textbf{Selected task scope:} Hermitian nonnegative sesquilinear S on complex polynomials. It is the spectral function of Theta=\allowbreak{}(J3,J5,alpha,beta), with J3 real Jacobi positive subdiagonal,J5 real symmetric five-diagonal positive second subdiagonal,alpha>0,beta real, iff S(u,v)=\allowbreak{}integral u(A)(1) conjugate(v(A)(1)) d sigma, for a probability measure sigma on R with all finite moments and strictly positive integral |g|\textasciicircum{}2 for every nonzero polynomial. A has exactly polynomial domain, raises x\textasciicircum{}k to degree k+\allowbreak{}1 with positive real leading coefficient and real remaining coefficients, and A Lambda0 is symmetric on polynomials. No boundedness, self-adjointness or symmetry of A itself is assumed.\par
\noindent\textbf{Difficulty:} H2: The inverse construction must recover a real symmetric five-diagonal matrix from a generally non-symmetric polynomial-domain operator, while controlling the positivity of its second subdiagonal. It simultaneously reconstructs the two initial pencil parameters and proves equality of the model operator on both the constant and positive-degree polynomial sectors. This generalized inverse spectral problem exceeds the ordinary three-term Favard correspondence.\par
\noindent\textbf{Editorial boundary note (not author text):} The complete original Theorem 3.1, both necessity and sufficiency, and full new triangular associated-operator Lemma 3.2 are retained. The exact Hermitian/nonnegative spectral-functional conventions, polynomial-domain integral model, normalized initial p1/alpha/beta reconstruction, real symmetric five-diagonal recovery and positive second subdiagonal are supplied. The reconstruction identifies the model operator both on x times polynomials and on the constant. Favard and the prior associated-polynomial model are explicitly cited established background. No symmetry/self-adjointness/boundedness of A itself, unique integral measure, multiplication-operator simplification or later perturbation formula is asserted; no neighboring corollary is counted.\par
\noindent\textbf{Source:} \url{https://sigma-journal.com/2017/085/}\quad DOI: \url{https://doi.org/10.3842/SIGMA.2017.085}\par
\noindent\textbf{License and attribution:} Copyright retained by the named authors. Published by SIGMA under CC BY 4.0: \url{https://creativecommons.org/licenses/by/4.0/}. Source: \url{https://sigma-journal.com/about.html}. Adaptation consists of standalone excerpt formatting, cover and numbering restoration; original author language and mathematical text are retained.\par
\medskip\hrule\medskip\noindent\textbf{Author text begins below.}\par
\setcounter{section}{1}
\setcounter{subsection}{0}
\setcounter{subsubsection}{0}
\setcounter{Theorem}{1}
\setcounter{equation}{1}
% Original source lines 85--121
\begin{Definition}\label{d1_1}A set $\Theta = (J_3, J_5, \alpha, \beta)$, where $\alpha>0$, $\beta\in\mathbb{R}$, $J_3$ is a Jacobi matrix and~$J_5$ is a semi-inf\/inite real symmetric f\/ive-diagonal matrix with positive numbers on the second subdiagonal, is said to be \textit{a Jacobi-type pencil $($of matrices$)$}.
\end{Definition}
From this def\/inition we see that matrices $J_3$ and $J_5$ have the following form
\begin{gather}
J_3 =
\left(
\begin{matrix}
b_0 & a_0 & 0 & 0 & 0 & \cdots\\
a_0 & b_1 & a_1 & 0 & 0 & \cdots\\
0 & a_1 & b_2 & a_2 & 0 &\cdots\\
\vdots & \vdots & \vdots & \ddots \end{matrix}
\right),\qquad a_k>0,\quad b_k\in\mathbb{R},\quad k\in\mathbb{Z}_+, \label{f1_5} \\
J_5 = \left(
\begin{matrix}
\alpha_0 & \beta_0 & \gamma_0 & 0 & 0 & 0 & \cdots\\
\beta_0 & \alpha_1 & \beta_1 & \gamma_1 & 0 & 0 & \cdots\\
\gamma_0 & \beta_1 & \alpha_2 & \beta_2 & \gamma_2 & 0 & \cdots\\
0 & \gamma_1 & \beta_2 & \alpha_3 & \beta_3 & \gamma_3 & \cdots \\
\vdots & \vdots & \vdots &\vdots & \ddots \end{matrix}
\right),\qquad \alpha_n,\beta_n\in\mathbb{R},\quad \gamma_n>0,\quad n\in\mathbb{Z}_+. \label{f1_10}
\end{gather}

With a Jacobi-type pencil of matrices $\Theta$ one associates a system of polynomials $\{ p_n(\lambda) \}_{n=0}^\infty$, such that
\begin{gather*} p_0(\lambda) = 1,\qquad p_1(\lambda) = \alpha\lambda + \beta, \end{gather*}
and
\begin{gather}\label{f1_20}
(J_5 - \lambda J_3) \vec p(\lambda) = 0,
\end{gather}
where $\vec p(\lambda) = (p_0(\lambda), p_1(\lambda), p_2(\lambda),\dots)^{\rm T}$. Here the superscript ${\rm T}$ means the transposition. Polynomials $\{ p_n(\lambda) \}_{n=0}^\infty$ are said to be \textit{associated to the Jacobi-type pencil of matrices~$\Theta$}.

Observe that for each system of orthonormal polynomials on the real line with $p_0=1$ one may choose $J_3$ to be the corresponding Jacobi matrix (which elements are the recurrence coef\/f\/icients), $J_5 = J_3^2$, and $\alpha$, $\beta$ being the coef\/f\/icients of~$p_1$ ($p_1(\lambda) = \alpha\lambda + \beta$).

One can rewrite relation~(\ref{f1_20}) in the scalar form
\begin{gather} \gamma_{n-2} p_{n-2}(\lambda) + (\beta_{n-1}-\lambda a_{n-1}) p_{n-1}(\lambda) + (\alpha_n-\lambda b_n) p_n(\lambda) \notag \\
\qquad{} + (\beta_n-\lambda a_n) p_{n+1}(\lambda) + \gamma_n p_{n+2}(\lambda) = 0,\qquad n\in\mathbb{Z}_+, \label{f1_30}
\end{gather}
where $p_{-2}(\lambda) = p_{-1}(\lambda) = 0$, $\gamma_{-2} = \gamma_{-1} = a_{-1} = \beta_{-1} = 0$.

\setcounter{section}{1}
\setcounter{subsection}{0}
\setcounter{subsubsection}{0}
\setcounter{Theorem}{2}
\setcounter{equation}{5}
% Original source lines 141--399
{\bf Notations.} As usual, we denote by $\mathbb{R}$, $\mathbb{C}$, $\mathbb{N}$, $\mathbb{Z}$, $\mathbb{Z}_+$, the sets of real numbers, complex numbers, positive integers, integers and non-negative integers, respectively. By~$\mathbb{P}$ we denote the set of all polynomials with complex coef\/f\/icients.

By $l_2$ we denote the usual Hilbert space of all complex sequences $c = (c_n)_{n=0}^\infty = (c_0,c_1,c_2,\dots)^{\rm T}$ with the f\/inite norm $\| c \|_{l_2} = \sqrt{\sum\limits_{n=0}^\infty |c_n|^2}$. Here $T$ means the transposition. The scalar product of two sequences $c = (c_n)_{n=0}^\infty$, $d = (d_n)_{n=0}^\infty\in l_2$ is given by $(c,d)_{l_2} = \sum\limits_{n=0}^\infty c_n \overline{d_n}$. We denote $\vec e_m = (\delta_{n,m})_{n=0}^\infty\in l_2$, $m\in\mathbb{Z}_+$. By $l_{2,{\rm f\/in}}$ we denote the set of all f\/inite vectors from~$l_2$, i.e., vectors with all, but f\/inite number, components being zeros. By $\operatorname{diag} (c_1,c_2,c_3,\dots)$ we mean a~semi-inf\/inite matrix with $c_j\in\mathbb{C}$ in $j$-th column and $j$-th row and zeros outside the main diagonal.

By $\mathfrak{B}(\mathbb{R})$ we denote the set of all Borel subsets of $\mathbb{R}$. If $\sigma$ is a (non-negative) bounded measure on~$\mathfrak{B}(\mathbb{R})$ then by $L^2_\sigma$ we denote a Hilbert space of all (classes of equivalences of) complex-valued functions $f$ on $\mathbb{R}$ with a f\/inite norm $\| f \|_{L^2_\sigma} = \sqrt{ \int_\mathbb{R} |f(x)|^2 {\rm d}\sigma }$. The scalar product of two functions $f,g\in L^2_\sigma$ is given by $(f,g)_{L^2_\sigma} = \int_{\mathbb{R}} f(x) \overline{g(x)} {\rm d}\sigma$. By $[ f ]$ we denote the class of equivalence in~$L^2_\sigma$ which contains the representative~$f$. By $\mathcal{P}$ we denote a set of all (classes of equivalence which contain) polynomials in $L^2_\sigma$. As usual, we sometimes use the representatives instead of their classes in formulas. Let $B$ be an arbitrary linear operator in $L^2_\sigma$ with the domain $\mathcal{P}$. Let $f(\lambda)\in\mathbb{P}$ be nonzero and of degree $d\in\mathbb{Z}_+$, $f(\lambda) = \sum\limits_{k=0}^d d_k \lambda^k$, $d_k\in\mathbb{C}$. We set
\begin{gather*} f(B) = \sum_{k=0}^d d_k B^k,\qquad B^0 := E|_{\mathcal{P}}. \end{gather*}
If $f\equiv 0$, then $f(B):= 0|_{\mathcal{P}}$.

If H is a Hilbert space then $(\cdot,\cdot)_H$ and $\| \cdot \|_H$ mean the scalar product and the norm in $H$, respectively. Indices may be omitted in obvious cases. For a linear operator $A$ in $H$, we denote by $D(A)$ its domain, by $R(A)$ its range, by $\operatorname{Ker} A$ its null subspace (kernel), and $A^*$ means the adjoint operator if it exists. If $A$ is invertible then $A^{-1}$ means its inverse. $\overline{A}$ means the closure of the operator, if the operator is closable. If $A$ is bounded then $\| A \|$ denotes its norm. For a~set $M\subseteq H$ we denote by $\overline{M}$ the closure of $M$ in the norm of $H$. By $\operatorname{Lin} M$ we mean the set of all linear combinations of elements of $M$, and $\overline{\operatorname{span}}\, M := \overline{\operatorname{Lin} M}$. By $E=E_H$ ($0=0_H$) we denote the identity operator in $H$, i.e., $E_H x = x$, $x\in H$ (respectively the null operator in $H$, i.e., $0_H x = 0$, $x\in H$). If $H_1$ is a subspace of $H$, then $P_{H_1} = P_{H_1}^{H}$ is an operator of the orthogonal projection on $H_1$ in~$H$.

\section{Preliminaries}

In this section, for the convenience of the reader, we recall basic def\/initions and results from~\cite{cit_95000_Z}. Then we state the direct and inverse spectral problems for the Jacobi-type pencils. The direct problem will be solved immediately, while the inverse problem will be considered in the next section. Set
\begin{gather*}
u_n := J_3 \vec e_n = a_{n-1} \vec e_{n-1} + b_n \vec e_n + a_n \vec e_{n+1}, \\
w_n := J_5 \vec e_n = \gamma_{n-2} \vec e_{n-2} + \beta_{n-1} \vec e_{n-1} + \alpha_n \vec e_n + \beta_n \vec e_{n+1}+ \gamma_n \vec e_{n+2},\qquad n\in\mathbb{Z}_+.
\end{gather*}
Here and in what follows by $\vec e_k$ with negative $k$ we mean (vector) zero. The following operator
\begin{gather} A f = \frac{\zeta}{\alpha} (\vec e_1 - \beta \vec e_0) + \sum_{n=0}^\infty \xi_n w_n, \notag \\
f = \zeta \vec e_0 + \sum_{n=0}^\infty \xi_n u_n\in l_{2,{\rm f\/in}},\qquad \zeta, \xi_n\in\mathbb{C}, \label{f3_60}
\end{gather}
with $D(A) = l_{2,{\rm f\/in}}$ is said to be \textit{the associated operator for the Jacobi-type pencil $\Theta$}. Notice that in the sums in~(\ref{f3_60}) only f\/inite number of $\xi_n$ are nonzero. We shall always assume this in the case of elements from the linear span. In particular, we have
\begin{gather*} A J_3 \vec e_n = J_5 \vec e_n,\qquad n\in\mathbb{Z}_+, \end{gather*}
and therefore
\begin{gather*} A J_3 = J_5. \end{gather*}
As usual, the matrices $J_3$ and $J_5$ def\/ine linear operators with the domain $l_{2,{\rm f\/in}}$ which we denote by the same letters.

For an arbitrary non-zero polynomial $f(\lambda)\in\mathbb{P}$ of degree $d\in\mathbb{Z}_+$, $f(\lambda) = \sum\limits_{k=0}^d d_k \lambda^k$, $d_k\in\mathbb{C}$, we set $f(A) = \sum\limits_{k=0}^d d_k A^k$. Here $A^0 := E|_{l_{2,{\rm f\/in}}}$. For $f(\lambda) \equiv 0$, we set $f(A) = 0|_{l_{2,{\rm f\/in}}}$. The following relations hold~\cite{cit_95000_Z}
\begin{gather}\label{f3_110}
\vec e_n = p_n(A) \vec e_0,\qquad n\in\mathbb{Z}_+,\\
\label{f3_120} \big( p_n(A) \vec e_0, p_m(A) \vec e_0 \big)_{l_2} = \delta_{n,m},\qquad n,m\in\mathbb{Z}_+.
\end{gather}
Denote by $\{ r_n(\lambda) \}_{n=0}^\infty$, $r_0(\lambda) = 1$, the system of polynomials satisfying
\begin{gather}
\label{f3_130}
J_3 \vec r(\lambda) = \lambda \vec r(\lambda),\qquad \vec r(\lambda) = (r_0(\lambda), r_1(\lambda),r_2(\lambda), \dots )^{\rm T}.
\end{gather}
These polynomials are orthonormal on the real line with respect to a (possibly non-unique) non-negative f\/inite measure $\sigma$ on~$\mathfrak{B}(\mathbb{R})$ (Favard's theorem). Consider the following operator
\begin{gather}\label{f3_140}
U \sum_{n=0}^\infty \xi_n \vec e_n = \left[ \sum_{n=0}^\infty \xi_n r_n(x) \right ],\qquad \xi_n\in\mathbb{C},
\end{gather}
which maps $l_{2,{\rm f\/in}}$ onto $\mathcal{P}$. Here by $\mathcal{P}$ we denote a set of all (classes of equivalence which contain) polynomials in $L^2_\sigma$. Denote
\begin{gather} \label{f3_150}
\mathcal{A} = \mathcal{A}_\sigma = U A U^{-1}.
\end{gather}
The operator $\mathcal{A} = \mathcal{A}_\sigma$ is said to be \textit{the model representation in $L^2_\sigma$ of the associated operator $A$}.

\begin{Theorem}[\cite{cit_95000_Z}]\label{t3_1} Let $\Theta = (J_3, J_5, \alpha, \beta)$ be a Jacobi-type pencil. Let $\{ r_n(\lambda) \}_{n=0}^\infty$, $r_0(\lambda) = 1$, be a system of polynomials satisfying~\eqref{f3_130} and $\sigma$ be their $($arbitrary$)$ orthogonality measure on~$\mathfrak{B}(\mathbb{R})$. The associated polynomials $\{ p_n(\lambda) \}_{n=0}^\infty$ satisfy the following relations
\begin{gather}\label{f3_160}
\int_{\mathbb{R}} p_n(\mathcal{A})(1) \overline{ p_m(\mathcal{A})(1) } {\rm d}\sigma = \delta_{n,m},\qquad n,m\in\mathbb{Z}_+,
\end{gather}
where $\mathcal{A}$ is the model representation in $L^2_\sigma$ of the associated operator $A$.
\end{Theorem}

We have now recalled all basic notations and results which we shall need in the present paper. We introduce the following def\/inition.
\begin{Definition} \label{dd1_5} Let $\Theta = (J_3, J_5, \alpha, \beta)$ be a Jacobi-type pencil and $\{ p_n(\lambda) \}_{n=0}^\infty$ be the associated polynomials to $\Theta$. A sesquilinear functional $S(u,v)$, $u,v\in\mathbb{P}$, satisfying the following relation
\begin{gather*} S(p_n,p_m) = \delta_{n,m},\qquad n,m\in\mathbb{Z}_+, \end{gather*}
is said to be \textit{the spectral function of the Jacobi-type pencil $\Theta$}.
\end{Definition}

In a usual manner, we may introduce the corresponding direct and inverse spectral problems. \textit{The direct spectral problem for a Jacobi-type pencil~$\Theta$} consists in searching for answers on the following questions:
\begin{itemize}\itemsep=0pt
 \item[1)] Does the spectral function exist?
 \item[2)] If the spectral function exists, is it unique?
 \item[3)] If the spectral function exists, how to f\/ind it (or them)?
\end{itemize}

By relation~(\ref{f3_120}) we see that the spectral function always exists. By the linearity it follows that the spectral function is unique. It has the following representation
\begin{gather} \label{ff1_7}
S(u,v) = \big( u(A) \vec e_0, v(A) \vec e_0 \big)_{l_2},\qquad u,v\in\mathbb{P},
\end{gather}
where $A$ is the associated operator for $\Theta$. By~(\ref{ff1_7}) we obtain that
\begin{gather}\label{ff1_7_0}
\overline{S(u,v)} = S(v,u),\qquad u,v\in\mathbb{P},
\end{gather}
and
\begin{gather} \label{ff1_7_0_1}
S(u,u)\geq 0,\qquad u\in\mathbb{P}.
\end{gather}
Moreover, the following integral representation holds
\begin{gather} \label{ff1_7_1}
S(u,v) = \int_{\mathbb{R}} u(\mathcal{A}_\sigma)(1) \overline{ v(\mathcal{A}_\sigma)(1) } {\rm d}\sigma,\qquad u,v\in\mathbb{P}.
\end{gather}
Thus, the direct spectral problem for $\Theta$ is solved in full.

\textit{The inverse spectral problem for a Jacobi-type pencil $\Theta$} consists in searching for answers on the following questions:
\begin{itemize}\itemsep=0pt
\item[(a)] Is it possible to reconstruct the Jacobi-type pencil $\Theta = (J_3, J_5, \alpha, \beta)$ using its spectral function? If it is possible, what is the procedure of the reconstruction?
\item[(b)] What are necessary and suf\/f\/icient conditions for a sesquilinear functional $\sigma(u,v)$, $u,v\in\mathbb{P}$, to be the spectral function of a Jacobi-type pencil $\Theta = (J_3, J_5, \alpha, \beta)$?
\end{itemize}
Questions~(a) and (b) will be investigated in the next section.

\section{The inverse spectral problem for a Jacobi-type pencil}

An answer to the question~(b) of the inverse spectral problem is given by the following theorem.

\begin{Theorem}\label{tt2_1}A sesquilinear functional $S(u,v)$, $u,v\in\mathbb{P}$, satisfying relations~\eqref{ff1_7_0},~\eqref{ff1_7_0_1}, is the spectral function of a Jacobi-type pencil if and only if it admits the following integral representation
\begin{gather}\label{ff2_5}
S(u,v) = \int_{\mathbb{R}} u(\mathcal{A})(1) \overline{ v(\mathcal{A})(1) } {\rm d}\sigma,\qquad u,v\in\mathbb{P},
\end{gather}
where $\sigma$ is a non-negative measure on $\mathfrak{B}(\mathbb{R})$ with all finite power moments,
\begin{gather*} \int_\mathbb{R} {\rm d}\sigma = 1, \qquad \int_\mathbb{R} |g(x)|^2 {\rm d}\sigma > 0,
\end{gather*} for any non-zero complex polynomial~$g$, and $\mathcal{A}$ is a linear operator in~$L^2_\sigma$ with the following properties:
\begin{itemize}\itemsep=0pt
\item[$(i)$] $D(\mathcal{A}) = \mathcal{P}$;
\item[$(ii)$] for each $k\in\mathbb{Z}_+$ it holds
\begin{gather}\label{ff2_7}
\mathcal{A} x^k = \xi_{k,k+1} x^{k+1} + \sum_{j=0}^k \xi_{k,j} x^j,
\end{gather}
where $\xi_{k,k+1} > 0$, $\xi_{k,j}\in\mathbb{R}$, $0\leq j\leq k$;
\item[$(iii)$] the operator $\mathcal{A} \Lambda_0$ is symmetric, here by~$\Lambda_0$ we denote the operator of the multiplication by an independent variable in~$L^2_\sigma$ restricted to~$\mathcal{P}$.
\end{itemize}
\end{Theorem}
\begin{proof} \textit{Necessity.} Let $\Theta = ( J_3, J_5, \alpha, \beta )$ be a Jacobi-type pencil and $A$ be the associated operator of~$\Theta$. Def\/ine polynomials $\{ r_n(\lambda) \}$, the measure $\sigma$, the operators $U$ and $\mathcal{A}=\mathcal{A}_\sigma$ as in the introduction, see~(\ref{f3_130}), (\ref{f3_140}), (\ref{f3_150}). Let $S(u,v)$ be the spectral function of the pencil~$\Theta$. It has an integral representation~(\ref{ff1_7_1}). It remains to check that $\sigma$ and $\mathcal{A}$ possess all the required properties in the statement of the theorem.

Since $\sigma$ is generated by the Jacobi matrix $J_3$, it has all required properties as in the statement of the theorem. The operator~$\mathcal{A}_\sigma$ is linear and def\/ined on~$\mathcal{P}$.

\begin{Lemma}\label{ll2_1} The associated operator $A$ of a Jacobi-type pencil $\Theta$ has the following properties:
\begin{itemize}\itemsep=0pt
\item[$1)$] for each $n\in\mathbb{Z}_+$ it holds
\begin{gather} \label{ff2_10}
\vec e_n = d_{n,n} A^n \vec e_0 + \sum_{j=0}^{n-1} d_{n,j} A^j \vec e_0,
\end{gather}
where $d_{n,n} > 0$, $d_{n,j}\in\mathbb{R}$, $0\leq j\leq n-1$;

\item[$2)$] for each $n\in\mathbb{Z}_+$ it holds
\begin{gather}\label{ff2_12}
A^n \vec e_0 = c_{n,n} \vec e_n + \sum_{j=0}^{n-1} c_{n,j} \vec e_j,
\end{gather}
where $c_{n,n} > 0$, $c_{n,j}\in\mathbb{R}$, $0\leq j\leq n-1$;

\item[$3)$] for each $n\in\mathbb{Z}_+$ it holds
\begin{gather}\label{ff2_14}
A \vec e_n = f_{n,n+1} \vec e_{n+1} + \sum_{j=0}^{n} f_{n,j} \vec e_j,
\end{gather}
where $f_{n,n+1} > 0$, $f_{n,j}\in\mathbb{R}$, $0\leq j\leq n$.
\end{itemize}
\end{Lemma}

\begin{proof}[Proof of Lemma~\ref{ll2_1}] Let $p_n(\lambda) = \sum\limits_{j=0}^n d_{n,j} \lambda^j$, $d_{n,n}>0$, $d_{n,j}\in\mathbb{R}$, be the associated polyno\-mials of the pencil. By~(\ref{f3_110}) we conclude that relation~(\ref{ff2_10}) holds.

Let us check relation~(\ref{ff2_12}) by the induction argument. For the cases $n=0,1$ relation~(\ref{ff2_12}) obviously holds, see the def\/inition of $A$. Assume that~(\ref{ff2_12}) holds for $n=1,2,\dots ,r$, $r\in\mathbb{N}$. By~(\ref{ff2_10}) with $n=r+1$ we may write
\begin{gather*} A^{r+1} \vec e_0 = \frac{1}{d_{r+1,r+1}} \vec e_{r+1} - \frac{1}{d_{r+1,r+1}} \sum_{j=0}^{r} d_{r+1,j} A^j \vec e_0 \\
\hphantom{A^{r+1} \vec e_0}{}
= \frac{1}{d_{r+1,r+1}} \vec e_{r+1} + \sum_{j=0}^{r} \sum_{m=0}^{j} \frac{ (-1) d_{r+1,j} c_{j,m} }{d_{r+1,r+1}}\vec e_m.
\end{gather*}
Finally, in order to prove relation~(\ref{ff2_14}) we apply the operator $A$ to the both sides of relation~(\ref{ff2_10}) and use~(\ref{ff2_12})
\begin{gather*} A \vec e_n = d_{n,n} A^{n+1} \vec e_0 + \sum_{j=0}^{n-1} d_{n,j} A^{j+1} \vec e_0 \\
\hphantom{A \vec e_n}{}
= d_{n,n} c_{n+1,n+1} \vec e_{n+1} + \sum_{j=0}^{n} d_{n,n} c_{n+1,j} \vec e_j + \sum_{j=0}^{n-1} \sum_{k=0}^{j+1} d_{n,j} c_{j+1,k} \vec e_k,\qquad n\in\mathbb{Z}_+.
\end{gather*}
Lemma~\ref{ll2_1} is proved.
\end{proof}

Return to the proof of Theorem~\ref{tt2_1}. Choose an arbitrary $k\in\mathbb{Z}_+$. Observe that
\begin{gather*} x^k = \sum_{j=0}^k g_{k,j} r_j(x), \qquad g_{k,j}\in\mathbb{R},\quad g_{k,k}>0. \end{gather*}
By~(\ref{ff2_14}) we may write
\begin{gather} \mathcal{A}_\sigma \big[ x^k \big] = \sum_{j=0}^k g_{k,j} U A \vec e_j = \sum_{j=0}^k g_{k,j} \left(
f_{j,j+1} [r_{j+1}] + \sum_{l=0}^j f_{j,l} [r_l] \right) \notag \\
\hphantom{\mathcal{A}_\sigma \big[ x^k \big]}{} = g_{k,k} f_{k,k+1} [r_{k+1}] + \left[
\sum_{j=0}^{k-1} g_{k,j} f_{j,j+1} r_{j+1} + \sum_{j=0}^k g_{k,j} \sum_{l=0}^j f_{j,l} r_l\right] \notag \\
\hphantom{\mathcal{A}_\sigma \big[ x^k \big]}{}
= \left[ g_{k,k} f_{k,k+1} \widetilde{\mu_{k+1,k+1}} x^{k+1} + g_{k,k} f_{k,k+1} \sum_{m=0}^k \widetilde{\mu_{k+1,m}} x^{m} \right. \notag \\
\left.\hphantom{\mathcal{A}_\sigma \big[ x^k \big]=}{} +
\sum_{j=0}^{k-1} g_{k,j} f_{j,j+1} r_{j+1} + \sum_{j=0}^k g_{k,j} \sum_{l=0}^j f_{j,l} r_l \right], \label{ff2_15}
\end{gather}
where we set $r_l(\lambda) = \sum\limits_{j=0}^l \widetilde \mu_{l,j} \lambda^j$, $\widetilde\mu_{l,l}>0$, $\mu_{l,j}\in\mathbb{R}$. Since we get a~real polynomial of degree $k+1$ on the right-hand side of~(\ref{ff2_15}) and it has a positive leading coef\/f\/icient $g_{k,k} f_{k,k+1} \widetilde{\mu_{k+1,k+1}}$, then relation~(\ref{ff2_7}) is proved.

It remains to verify~$(iii)$. For arbitrary $n,m\in\mathbb{Z}_+$ we may write
\begin{gather}
(\mathcal{A}_\sigma\Lambda_0 [r_n(x)], [r_m(x)]) = \big(UAU^{-1} [a_{n-1} r_{n-1}(x) + b_n r_{n}(x) + a_n r_{n+1}(x)] , [r_m(x)]\big) \notag \\
\hphantom{(\mathcal{A}_\sigma\Lambda_0 [r_n(x)], [r_m(x)])}{}
 = (A J_3 \vec e_n, \vec e_m) = (J_5 \vec e_n, \vec e_m) = (\vec e_n, J_5 \vec e_m) \notag \\
\hphantom{(\mathcal{A}_\sigma\Lambda_0 [r_n(x)], [r_m(x)])}{}
= (\vec e_n, A (a_{m-1} \vec e_{m-1} + b_m \vec e_m + a_m \vec e_{m+1})) \nonumber\\
\hphantom{(\mathcal{A}_\sigma\Lambda_0 [r_n(x)], [r_m(x)])}{} = ([r_n(x)], \mathcal{A}_\sigma\Lambda_0 [r_m(x)]). \label{ff2_16}
\end{gather}
By the linearity we conclude that $\mathcal{A}_\sigma\Lambda_0$ is symmetric.

\textit{Sufficiency.} Suppose that a sesquilinear functional $S(u,v)$, $u,v\in\mathbb{P}$, satisfying relations~(\ref{ff1_7_0}), (\ref{ff1_7_0_1}), is given. Assume that it has the integral representation~(\ref{ff2_5}) where $\sigma$ is a non-negative probability measure on~$\mathfrak{B}(\mathbb{R})$ with all f\/inite moments, $\int_\mathbb{R} |g(x)|^2 {\rm d}\sigma > 0$, for any non-zero complex polynomial~$g$, and $\mathcal{A}$ is a linear operator in $L^2_\sigma$ with properties~$(i)$--$(iii)$. By condition~$(ii)$ and the induction argument it can be directly verif\/ied that for each $n\in\mathbb{Z}_+$ it holds
\begin{gather*} \mathcal{A}^n [1] = a_{n,n} [x^n] + \sum_{j=0}^{n-1} a_{n,j} \big[x^j\big], \end{gather*}
where $a_{n,n} > 0$, $a_{n,j}\in\mathbb{R}$, $0\leq j\leq n-1$.

Suppose that $S(u,u)=0$ for a complex polynomial $u$. Then
\begin{gather*} 0 = \| u(\mathcal{A})(1) \|^2_{L^2_\sigma}. \end{gather*}
Assume that $u$ is nonzero, $u(\lambda) = \sum\limits_{k=0}^n d_k \lambda^k$, $d_n\not=0$, $d_k\in\mathbb{C}$, $n\in\mathbb{Z}_+$. Observe that
\begin{gather*}
u(\mathcal{A})(1) = \left( \sum_{k=0}^n d_k \mathcal{A}^k \right) [1] = d_n \mathcal{A}^n [1] + \sum_{k=0}^{n-1} d_k \mathcal{A}^k [1] \\
\hphantom{u(\mathcal{A})(1)}{} = d_n \left( a_{n,n} \big[x^n\big] + \left[ \sum_{j=0}^{n-1} a_{n,j} x^j\right]\right)
+ \sum_{k=0}^{n-1} d_k \sum_{j=0}^k a_{k,j} \big[x^j\big] \\
\hphantom{u(\mathcal{A})(1)}{}
= \left[ d_n a_{n,n} x^n + d_n \sum_{j=0}^{n-1} a_{n,j} x^j + \sum_{k=0}^{n-1} \sum_{j=0}^k d_k a_{k,j} x^j \right] =: [r(x)].
\end{gather*}
We have a nonzero polynomial $r$ (of degree $n$) with $\| r \|_{L^2_\sigma} = 0$. This contradicts to our assumptions on the measure $\sigma$. Consequently, $u\equiv 0$.

The functional $S$ def\/ines an inner product on~$\mathbb{P}$. Thus, the complex vector space $\mathbb{P}$ becomes a space $H$ with a scalar product. It is a normed space with the norm $\| p \| = \sqrt{S(p,p)}$. We shall not need its completion. Set
\begin{gather*} \mathbf{p}_0(\lambda) = \frac{1}{ \| 1 \|_{H} } =1,\qquad
 \mathbf{p}_1(\lambda) = \frac{\lambda - S(x,1) 1}{ \| \lambda - S(x,1) 1 \|_{H} }. \end{gather*}
Let us check that $\mathbf{p}_1(\lambda)$ is well-def\/ined. Denote by $\{ s_k \}_{k=0}^\infty$ the power moments of $\sigma$
\begin{gather*} s_k = \int_\mathbb{R} x^k {\rm d}\sigma,\qquad k\in\mathbb{Z}_+. \end{gather*}
Let
\begin{gather*} \Delta_n := \det (s_{k+l})_{k,l=0}^n,\qquad n\in\mathbb{Z}_+,\qquad \Delta_{-1}:=1, \end{gather*}
be the corresponding Hankel determinants. Observe that
\begin{gather*} S(x,1) = \int \mathcal{A} (1) {\rm d}\sigma = \int (\xi_{0,1}\lambda + \xi_{0,0}) {\rm d}\sigma
= \xi_{0,1} s_1 + \xi_{0,0}\in\mathbb{R}, \\
 S(x,x) = \int \mathcal{A} (1) \overline{\mathcal{A} (1)} {\rm d}\sigma = \int (\xi_{0,1}\lambda + \xi_{0,0})^2 {\rm d}\sigma = \xi_{0,1}^2 s_2 + 2\xi_{0,1}\xi_{0,0} s_1 + \xi_{0,0}^2.
\end{gather*}
We may write
\begin{gather*} \| \lambda - S(x,1) 1 \|_{H}^2 = S(\lambda - S(x,1) 1, \lambda - S(x,1) 1) =
S(\lambda,\lambda) - (S(\lambda,1))^2 \\
 \hphantom{\| \lambda - S(x,1) 1 \|_{H}^2}{} = \xi_{0,1}^2 \big(s_2 - s_1^2\big) = \xi_{0,1}^2 \Delta_1 > 0,
\end{gather*}
where we have used our assumptions on the measure $\sigma$. Then
\begin{gather*} \mathbf{p}_1(\lambda) = \frac{\lambda - \xi_{0,1} s_1 - \xi_{0,0}}{ \xi_{0,1}\sqrt{\Delta_1} }. \end{gather*}
Set
\begin{gather}\label{ff2_20_5}
\alpha = \frac{1}{ \xi_{0,1}\sqrt{\Delta_1} },\qquad \beta = -\frac{\xi_{0,1} s_1 + \xi_{0,0}}{ \xi_{0,1}\sqrt{\Delta_1} }.
\end{gather}

Let $\{ \mathbf{r}_n(\lambda) \}_{n=0}^\infty$ be orthonormal polynomials with respect to the measure $\sigma$ (having positive leading coef\/f\/icients). Denote by $J_3$ the corresponding Jacobi matrix (formed by the recurrence coef\/f\/icients of $\mathbf{r}_n$). Denote
\begin{gather*} J_5 = (g_{m,n})_{m,n=0}^\infty,\qquad g_{m,n} := (\mathcal{A} \Lambda_0 [ \mathbf{r}_n(\lambda)], [ \mathbf{r}_m(\lambda)])_{L^2_\sigma}. \end{gather*}
By condition~$(iii)$ of the theorem we conclude that $J_5$ is a symmetric semi-inf\/inite matrix. Let
\begin{gather*} \mathbf{r}_n(\lambda) = \sum_{k=0}^n \eta_{n,k} \lambda^k,\qquad \eta_{n,k}\in\mathbb{R},\quad \eta_{n,n} > 0,\quad n\in\mathbb{Z}_+. \end{gather*}
For an arbitrary $n\in\mathbb{Z}_+$ by condition~$(ii)$ we may write
\begin{gather}\label{ff2_25}
\mathcal{A} \Lambda_0 [ \mathbf{r}_n(\lambda)] = \sum_{k=0}^n \eta_{n,k} \mathcal{A} \big[\lambda^{k+1}\big] = \big[
\eta_{n,n} \xi_{n+1,n+2} \lambda^{n+2} + d_{n+1}(\lambda) \big],
\end{gather}
where $d_{n+1}(\lambda)$ is a zero polynomial or a polynomial with real coef\/f\/icients, $\deg p\leq n+1$. Then
\begin{gather*} g_{m,n} = 0,\qquad m,n\in\mathbb{Z}_+, \quad m>n+2, \end{gather*}
and
\begin{gather*} g_{n+2,n} = \eta_{n,n} \xi_{n+1,n+2} \big(\big[\lambda^{n+2}\big], [r_{n+2}(\lambda)]\big) =
\frac{\eta_{n,n} \xi_{n+1,n+2}}{\eta_{n+2,n+2}} > 0,\qquad n\in\mathbb{Z}_+. \end{gather*}
By~(\ref{ff2_25}) we also see that $g_{m,n}$ are real numbers. We conclude that $J_5$ is real f\/ive-diagonal and it has positive numbers on the second sub-diagonal.

Consider a Jacobi-type pencil $\widetilde\Theta = (J_3, J_5, \alpha,\beta)$. Let $\{ p_n(\lambda) \}_{n=0}^\infty$ be the associated polynomials to $\widetilde\Theta$. For the pencil $\widetilde\Theta$ we def\/ine the standard objects from the introduction. The polynomials~$r_n(\lambda)$ (see~(\ref{f3_130})) coincide with $\mathbf{r}_n(\lambda)$. We may also choose the given $\sigma$ as the orthogonality measure (the orthogonality measure can be non-unique). The operators~$A$,~$U$ and~$\mathcal{A}_\sigma$ we def\/ine in the standard way, see~(\ref{f3_60}), (\ref{f3_140}), (\ref{f3_150}). We only can not use the brief notation~$\mathcal{A}$ for~$\mathcal{A}_\sigma$, since~$\mathcal{A}$ already denotes the operator in the integral representation of~$S$.

For the pencil $\widetilde\Theta$ we may apply our arguments in the proved Necessity. By relation~(\ref{ff2_16}) we may write
\begin{gather*} (\mathcal{A}_\sigma \Lambda_0 [r_n],[r_m])_{L^2_\sigma} = (J_5 \vec e_n, \vec e_m)_{l_2} = g_{m,n} = (\mathcal{A} \Lambda_0 [r_n],[r_m])_{L^2_\sigma},\qquad n,m\in\mathbb{Z}_+. \end{gather*}
Here the last equality follows by the def\/inition of the matrix $J_5$. Therefore
\begin{gather}\label{ff2_30}
\mathcal{A}_\sigma [ \lambda p(\lambda) ] = \mathcal{A} [ \lambda p(\lambda) ],\qquad \forall\, p\in\mathbb{P}.
\end{gather}
Notice that
\begin{gather}\label{ff2_32}
\mathcal{A}_\sigma [ 1 ] = U A \vec e_0 =\frac{1}{\alpha} U (\vec e_1 - \beta \vec e_0)=\left[ \frac{1}{\alpha} (r_1(\lambda) - \beta) \right].
\end{gather}
The orthonormal polynomial $r_1(\lambda)$ has the following form
\begin{gather}\label{ff2_34}
r_1(\lambda) = \frac{1}{\sqrt{\Delta_1}} (\lambda - s_1).
\end{gather}
By~(\ref{ff2_32}), (\ref{ff2_34}) and (\ref{ff2_20_5}) we conclude that
\begin{gather*} \mathcal{A}_\sigma [ 1 ] = [ \xi_{0,1} \lambda + \xi_{0,0} ]. \end{gather*}
On the other hand, by property $(ii)$ we have $\mathcal{A} [ 1 ] = [ \xi_{0,1} \lambda + \xi_{0,0} ]$. Therefore
\begin{gather}\label{ff2_40}
\mathcal{A}_\sigma [ 1 ] = \mathcal{A} [ 1 ].
\end{gather}
Relations~(\ref{ff2_30}) and~(\ref{ff2_40}) show that $\mathcal{A}_\sigma = \mathcal{A}$. Comparing relations~(\ref{ff1_7_1}) and~(\ref{ff2_5}) we see that the spectral function of $\widetilde\Theta$ coincides with~$S$.
\end{proof}
\begin{thebibliography}{99}
\bibitem[16]{cit_95000_Z}
Zagorodnyuk S.M., Orthogonal polynomials associated with some {J}acobi-type
 pencils, \href{https://doi.org/10.1007/s11253-017-1300-3}{\textit{Ukrain. Math.~J.}} \textbf{68} (2016), 1353--1365,
 \href{https://arxiv.org/abs/1508.01794}{arXiv:1508.01794}.

\end{thebibliography}
\end{document}
